\chapter{2D Truss and Beam Elements}
\label{ch:truss_beam}

\section{2D Truss Element}

A truss element can carry axial loads only. In a 2D structure, each node has
two translational DOFs: $u$ (horizontal) and $v$ (vertical). The element is
inclined at angle $\theta$ with respect to the global $x$-axis~\cite{cook2002concepts}.

\subsection{Local Stiffness Matrix}

In the local coordinate system aligned with the element axis, the stiffness
matrix is identical to the 1D bar element~\cref{eq:bar_stiffness}, expanded
to account for the transverse DOF (which has zero stiffness):
\begin{equation}
    \bar{\matr{k}}^e = \frac{EA}{L}
    \begin{bmatrix} 1 & 0 & -1 & 0 \\ 0 & 0 & 0 & 0 \\
                   -1 & 0 &  1 & 0 \\ 0 & 0 & 0 & 0 \end{bmatrix}
    \label{eq:truss_local}
\end{equation}
where the DOF ordering is $[\bar{u}_1,\, \bar{v}_1,\, \bar{u}_2,\, \bar{v}_2]$.

\subsection{Coordinate Transformation}

The transformation matrix $\matr{T}$ rotates from global $(x,y)$ to local
$(\bar{x}, \bar{y})$ coordinates. With $c = \cos\theta$, $s = \sin\theta$:
\begin{equation}
    \matr{T} = \begin{bmatrix}
         c & s & 0 & 0 \\
        -s & c & 0 & 0 \\
         0 & 0 & c & s \\
         0 & 0 &-s & c
    \end{bmatrix}
    \label{eq:truss_T}
\end{equation}

\begin{figure}[htbp]
\centering
\begin{tikzpicture}[>=Stealth, scale=1.15]
  %% Element from node 1 (0,0) to node 2 at 35 degrees, length ~4
  \coordinate (N1) at (0,0);
  \coordinate (N2) at ({4*cos(35)},{4*sin(35)});
  %% Bar
  \draw[line width=2.5pt,gray!55] (N1) -- (N2);
  \filldraw (N1) circle (3.5pt) node[below left=1pt] {\small 1};
  \filldraw (N2) circle (3.5pt) node[above right=1pt] {\small 2};
  %% Global axes
  \draw[->,thick] (N1) -- ++(2.2,0) node[right] {$x$};
  \draw[->,thick] (N1) -- ++(0,2.2) node[above] {$y$};
  %% Local axes (rotated by 35 deg)
  \draw[->,thick,blue!75!black] (N1) -- ++({2.0*cos(35)},{2.0*sin(35)}) node[right] {$\bar{x}$};
  \draw[->,thick,blue!75!black] (N1) -- ++({-2.0*sin(35)},{2.0*cos(35)}) node[above left] {$\bar{y}$};
  %% Angle arc and label
  \draw (1.1,0) arc[start angle=0, end angle=35, radius=1.1];
  \node at ({1.35*cos(17)},{1.35*sin(17)}) {$\theta$};
  %% DOF arrows at N2
  \draw[->,dashed,thick,red!65!black] (N2) -- ++(1.0,0) node[right] {\small $u_2$};
  \draw[->,dashed,thick,red!65!black] (N2) -- ++(0,1.0) node[above] {\small $v_2$};
  %% DOF label at N1
  \node[red!65!black,below right=2pt] at (N1) {\small $u_1,\,v_1$};
\end{tikzpicture}
\caption{2D truss element inclined at angle $\theta$: global axes $(x,y)$ and local axes $(\bar{x},\bar{y})$.}
\label{fig:truss_element}
\end{figure}

\subsection{Global Element Stiffness}

The element stiffness in the global frame is:
\begin{equation}
    \matr{k}^e = \matr{T}\transpose \bar{\matr{k}}^e \matr{T}
    = \frac{EA}{L}\begin{bmatrix}
         c^2  &  cs  & -c^2 & -cs  \\
         cs   &  s^2 & -cs  & -s^2 \\
        -c^2  & -cs  &  c^2 &  cs  \\
        -cs   & -s^2 &  cs  &  s^2
    \end{bmatrix}
    \label{eq:truss_global}
\end{equation}

\begin{importantbox}
\textbf{Angle convention:} $\theta$ is measured counter-clockwise from the
positive $x$-axis to the element axis (node 1 $\to$ node 2).
\end{importantbox}

\section{Euler-Bernoulli Beam Element}

The Euler-Bernoulli beam element models bending with the assumption that
plane cross-sections remain plane and perpendicular to the neutral axis
(no shear deformation).

\subsection{Degrees of Freedom}

Each node has two DOFs: transverse displacement $v$ and rotation $\theta$
(= $\dd v/\dd x$). The element has 4 DOFs total:
$\vect{u}^e = [v_1,\, \theta_1,\, v_2,\, \theta_2]\transpose$.

\subsection{Hermite Shape Functions}

Cubic Hermite polynomials ensure $C^1$ continuity (continuous displacement
and slope across elements):
\begin{align}
    H_1(x) &= 1 - 3\xi^2 + 2\xi^3 \notag \\
    H_2(x) &= L(\xi - 2\xi^2 + \xi^3) \notag \\
    H_3(x) &= 3\xi^2 - 2\xi^3 \notag \\
    H_4(x) &= L(-\xi^2 + \xi^3)
    \label{eq:hermite}
\end{align}
where $\xi = x/L \in [0,1]$.

\subsection{Beam Stiffness Matrix}

The curvature is $\kappa = \dd^2 v/\dd x^2 = \matr{B}_b\, \vect{u}^e$ and the
bending moment $M = EI\kappa$. The element stiffness is:
\begin{equation}
    \matr{k}^e_b = \int_0^L EI\, \matr{B}_b\transpose \matr{B}_b \dd x
    = \frac{EI}{L^3}
    \begin{bmatrix}
         12  &  6L  & -12  &  6L  \\
         6L  &  4L^2 & -6L  &  2L^2 \\
        -12  & -6L  &  12  & -6L  \\
         6L  &  2L^2 & -6L  &  4L^2
    \end{bmatrix}
    \label{eq:beam_stiffness}
\end{equation}

\begin{importantbox}
\begin{equation*}
    \matr{k}^e_b = \frac{EI}{L^3}
    \begin{bmatrix}
         12  &  6L  & -12  &  6L  \\
         6L  &  4L^2 & -6L  &  2L^2 \\
        -12  & -6L  &  12  & -6L  \\
         6L  &  2L^2 & -6L  &  4L^2
    \end{bmatrix}
\end{equation*}
\end{importantbox}

\subsection{Consistent Load Vector (Uniform Load $q$)}

\begin{equation}
    \vect{f}^e = \frac{qL}{12}
    \begin{bmatrix} 6 \\ L \\ 6 \\ -L \end{bmatrix}
    \label{eq:beam_load}
\end{equation}

\section{Frame Element}

A frame (beam-column) element combines axial bar stiffness and bending beam
stiffness. In the local frame the DOFs are:
$\vect{u}^e = [\bar{u}_1,\, \bar{v}_1,\, \bar{\theta}_1,\,
               \bar{u}_2,\, \bar{v}_2,\, \bar{\theta}_2]\transpose$
and the local stiffness is the $6\times6$ block-diagonal combination of
\cref{eq:bar_stiffness,eq:beam_stiffness}.

The $6\times6$ transformation matrix:
\begin{equation}
    \matr{T}_6 = \begin{bmatrix} c & s & 0 & 0 & 0 & 0 \\
                                 -s & c & 0 & 0 & 0 & 0 \\
                                  0 & 0 & 1 & 0 & 0 & 0 \\
                                  0 & 0 & 0 & c & s & 0 \\
                                  0 & 0 & 0 &-s & c & 0 \\
                                  0 & 0 & 0 & 0 & 0 & 1
                 \end{bmatrix}
    \label{eq:frame_T}
\end{equation}

The global frame element stiffness is $\matr{k}^e = \matr{T}_6\transpose
\bar{\matr{k}}^e_{\text{frame}} \matr{T}_6$.

\section{Worked Example: Simply Supported Beam}

\begin{example}[Two-element simply supported beam]
\label{ex:beam2}

A beam with $L = 2$ m, $EI = 10^4$ N$\cdot$m$^2$, simply supported at
both ends, with central point load $P = 1000$ N.

\medskip
\noindent Discretize with 2 equal beam elements ($L_e = 1$ m).
\medskip

\noindent\textbf{DOFs:} node 1: $(v_1, \theta_1)$, node 2: $(v_2, \theta_2)$,
node 3: $(v_3, \theta_3)$.

\noindent\textbf{BCs:} $v_1 = 0$, $v_3 = 0$ (pin supports).

\noindent\textbf{Free DOFs:} $[\theta_1, v_2, \theta_2, \theta_3]$.

\noindent\textbf{Analytical mid-span:} $v_2 = PL^3/(48EI) = 8.33 \times 10^{-4}$ m.

\noindent See \Cref{lst:beam2d} in the Appendix for the MATLAB implementation.
\end{example}
